\chapter{輸送計画としての量子チャネル}
\label{ch:quantum-channels}

有限集合上の Markov kernel は，入力分布を出力分布へ写す確率的な輸送計画である．
その非可換版が量子チャネルである．本章では Kraus 表現を定義として採用し，
完全正値性の一般論には立ち入らない．

\section{量子チャネル}

\begin{definition}{量子チャネル}{q-channel}
  有限次元 Hilbert 空間 $\mathcal{H},\mathcal{K}$ の間の\textbf{量子チャネル}とは，
  ある線形写像 $B_a\colon\mathcal{H}\to\mathcal{K}$ を用いて

  \[
    \Phi(A)=\sum_aB_aAB_a^*,
    \qquad
    \sum_aB_a^*B_a=\Id_{\mathcal{H}}
  \]

  と書ける線形写像 $\Phi\colon\Lin{\mathcal{H}}\to\Lin{\mathcal{K}}$ である．
  $(B_a)_a$ を Kraus 作用素という．
\end{definition}

量子チャネルは半正定値性と跡を保つので，量子状態を量子状態へ写す．

\begin{definition}{量子輸送計画}{q-transport-plan}
  $\rho\in\States{\mathcal{H}}$ から $\sigma\in\States{\mathcal{K}}$ への
  \textbf{量子輸送計画}とは

  \[
    \Phi(\rho)=\sigma
  \]

  をみたす量子チャネル $\Phi$ である．
\end{definition}

\begin{example}{置換チャネル}{q-replacement-channel}
  固定した $\sigma\in\States{\mathcal{K}}$ に対し

  \[
    \Phi_\sigma(A):=(\Tr A)\sigma
  \]

  は量子チャネルであり，すべての入力状態を $\sigma$ へ写す．したがって任意の二状態の間に
  量子輸送計画が存在する．これは入力を捨てて $\sigma$ を準備する操作である．
\end{example}

\begin{example}{Markov kernel の埋め込み}{q-markov-channel}
  有限集合 $\mathcal{X},\mathcal{Y}$ の間の Markov kernel $K(y\mid x)$ に対し

  \[
    B_{yx}:=\sqrt{K(y\mid x)}\ket{y}\bra{x}
  \]

  とおくと，$\sum_{x,y}B_{yx}^*B_{yx}=\Id$ である．対応する量子チャネルは，
  対角状態 $\rho_p$ を出力分布
  $q(y)=\sum_xK(y\mid x)p(x)$ の対角状態 $\rho_q$ へ写す．
\end{example}

\section{チャネルから coupling へ}

$\rho=\sum_jp_j\proj{e_j}$ とスペクトル分解し，双対空間 $\mathcal{H}^*$ を用いて

\[
  \ket{\sqrt\rho}
  :=\sum_j\sqrt{p_j}\,\ket{e_j}\otimes\bra{e_j}
  \in\mathcal{H}\otimes\mathcal{H}^*
\]

とおく．$\proj{\sqrt\rho}$ は純粋状態であり，その周辺状態は $\rho$ と転置 $\rho^T$ である．
これを $\rho$ の\textbf{標準純粋化}という．

\begin{proposition}{輸送計画が定める coupling}{q-channel-coupling}
  $\Phi(\rho)=\sigma$ をみたす量子チャネル $\Phi$ に対し

  \[
    \pi_\Phi
    :=(\Phi\otimes\Id)\bigl(\proj{\sqrt\rho}\bigr)
    \in\States{\mathcal{K}\otimes\mathcal{H}^*}
  \]

  とおくと

  \[
    \Tr_{\mathcal{H}^*}\pi_\Phi=\sigma,
    \qquad
    \Tr_{\mathcal{K}}\pi_\Phi=\rho^T
  \]

  が成り立つ．
\end{proposition}

\begin{proof}
  第1式は，$\mathcal{H}^*$ 側を部分跡した後に $\Phi$ を作用させることと，
  先に $\Phi$ を作用させてから部分跡することが一致することから
  $\Tr_{\mathcal{H}^*}\pi_\Phi=\Phi(\rho)=\sigma$ となる．
  第2式は，$\Phi$ が跡を保ち，$\mathcal{H}^*$ 側には作用しないことから従う．
\end{proof}

これは古典 coupling を条件付き分布，すなわち Markov kernel として読む操作の量子版である．
ただし，入力側に転置が現れ，純粋化が entanglement を含み得る点が古典系と異なる．

\section{二次コストの落とし穴}

観測量 $R\in\Obs{\mathcal{H}}$ から Euclid 距離の二乗を模した

\[
  C_R:=(R\otimes\Id-\Id\otimes R^T)^2
\]

を作り，チャネルから得られる coupling 上で平均コストを最小化することは自然に見える．
しかし，これは一般には距離を与えない．

\begin{example}{正の自己コスト}{q-positive-self-cost}
  $\mathcal{H}=\C^2$，$R=\proj{0}-\proj{1}$，
  $\ket{+}:=(\ket{0}+\ket{1})/\sqrt2$，$\rho=\proj{+}$ とする．
  一方の周辺状態が純粋な coupling は積状態に限るので，
  $\rho$ と $\rho$ の coupling は $\rho\otimes\rho^T$ だけである．ここで

  \[
    \bra{+}R\ket{+}=0,
    \qquad
    \bra{+}R^2\ket{+}=1
  \]

  だから

  \[
    \Tr\!\left[C_R(\rho\otimes\rho^T)\right]=2.
  \]

  したがって，この二次コストから作る量は同じ状態の間でも $0$ にならない．
\end{example}

そこで本資料の主目標には，coupling コストとは別の発想を使う．
多体系の「一つの qubit だけを変える」という局所性から単位球を作り，
その Minkowski 汎関数として量子 $W_1$ を定義する．
